Evaluation is framed around scientific utility and operational reliability rather than conversational quality alone.
This aligns with recent calls to evaluate agents across multi-step lifecycles and interactive analysis settings \citep{Zhang2024DSEval,Li2025IDABench}.

\paragraph{Current internal retrieval evaluation.}
A first evaluation has been performed at the level of the retrieval substrate, which provides the grounding evidence used both for question answering and for tool selection.
In an initial study ($N=108$) of reference-grounded queries whose answers are known to exist in the indexed corpus, the correct evidence passage is retrieved within the top-$32$ results in $95.4\%$ of cases, with median rank $1$.
These results should be interpreted as a preliminary characterization of retrieval robustness for technically specific facts (e.g., numeric thresholds and protocol details), rather than as an end-to-end assessment of the full agent.
A comprehensive evaluation that exercises the complete plan--execute--observe loop (including dataset/variable resolution, parameterization, and artifact production), curated and reviewed by domain scientists, is in preparation.

\paragraph{Benchmark prompt suite.}
A curated suite of prompts will cover: dataset discovery, variable resolution, spatiotemporal subsetting, visualization, climatology, and colocalization.
The suite can be seeded from maintained example prompts (included as supplementary material) and extended with domain-expert authored tasks.

\paragraph{Potential Metrics.}
\begin{itemize}[leftmargin=*]
  \item \textbf{Resolution correctness:} selection of the intended dataset table(s) and variable(s), judged against expert labels.
  \item \textbf{Tool success rate:} fraction of tasks completed without execution errors, including appropriate recovery on the first retry.
  \item \textbf{Output validity:} basic sanity checks on returned artifacts (e.g., non-empty subsets; expected dimensions; figure generation success).
  \item \textbf{Reproducibility:} ability to reproduce artifacts from emitted code or logged parameters.
  \item \textbf{Efficiency:} number of tool calls and end-to-end latency distributions (descriptive).
\end{itemize}

\paragraph{Baselines and comparisons.}
Comparisons may include: (i) direct use of the CMAP SDK without agent assistance, and (ii) ``chat-only'' LLM baselines without tool execution.
Where relevant, evaluation protocols can draw on broader agent benchmarks that emphasize interactive decision making \citep{Liu2023AgentBench} and tool-use performance \citep{Qin2023ToolLLM,Guo2024StableToolBench}.
